\documentclass[preprint,12pt]{elsarticle}

\usepackage[utf8]{inputenc}
\usepackage[T1]{fontenc}
\usepackage{graphicx}
\usepackage{booktabs}
\usepackage{amsmath,amssymb}
\usepackage{hyperref}
\usepackage{url}
\usepackage{lineno}

\journal{Online Social Networks and Media}

\begin{document}

\begin{frontmatter}

\title{Weaponizing Broadcast Channels: Command-and-Control Mobilization, Search Engine Manipulation, and Judicial AI Regulation in Brazil's 2026 Presidential Runoff}

\author[inst1]{André Luiz Satoshi Kawamoto\corref{cor1}}
\ead{kawamoto@utfpr.edu.br}
\cortext[cor1]{Corresponding author. Open Access publication supported under the CAPES--Elsevier Transformative Agreement (Brazil).}

\affiliation[inst1]{organization={Academic Department of Computing (DACOM), Federal University of Technology -- Paraná (UTFPR)},
            city={Campo Mourão},
            state={Paraná},
            country={Brazil}}

\begin{abstract}
Encrypted messaging applications such as WhatsApp have long been conceptualized in computational social science as decentralized networks of small, private conversational groups constrained by strict participant caps. However, the architectural rollout of \textit{WhatsApp Channels} fundamentally upended this paradigm by introducing one-to-many broadcast streams capable of reaching millions of followers with zero forwarding friction. In this paper, we conduct the first multi-method computational investigation into how rival presidential campaigns weaponized WhatsApp Channels as real-time Command-and-Control (C2) infrastructures during the 2026 Brazilian Presidential Election runoff between Luiz Inácio Lula da Silva and Flávio Bolsonaro. We present the theoretical framework, data triangulation pipeline, and formal modeling underpinning our longitudinal multi-platform study. Combining native message logs ($N=1,331$ messages across the complete 14-day observation window through October 30, 2026) from the two official channels registered with Brazil's Superior Electoral Court (TSE)---``Exército do Lula'' ($10\text{ million}$ followers) and ``Tropa do Flávio'' ($2.0\text{ million}$ followers)---with hourly Google Trends relative search volume, X (Twitter) Trending Topics scrapers, Wikimedia REST API pageviews, and unauthenticated HTTP subscriber tracking, our methodology enables empirical measurement across three core phenomena: (1) coordinated \textbf{Linkless Search Engine Manipulation (Google Bombing)}, where explicit text commands engineered nationwide search spikes while evading URL-filtering algorithms; (2) \textbf{cross-platform narrative brigading}, illustrating symmetric and asymmetric warfare across Instagram and X; and (3) \textbf{real-time campaign co-adaptation to judicial regulation}, capturing compliance with mandatory AI-content labeling and strategic gamification of pre-election shutdown threats. We formulate the methodological and conceptual architecture of the study, establishing structured analytical placeholders for the final 120-hour post-election empirical results.
\end{abstract}

\begin{highlights}
\item First comprehensive study of mega-scale WhatsApp Channels ($>10\text{M}$ followers) as political C2 infrastructures.
\item Formalizes the causal triangulation of linkless Google Search Bombing across Google Trends, X, and Wikipedia.
\item Documents tactical co-adaptation to Electoral Court (TSE) mandatory AI labeling and channel suspension threats.
\item Unpacks platform governance friction: default-muted broadcasts bypassed via notification social engineering.
\item Integrates longitudinal multi-platform telemetry across WhatsApp, Instagram, X, Google Search, and Wikipedia.
\end{highlights}

\begin{keyword}
WhatsApp Channels \sep Coordinated Inauthentic Behavior \sep Search Engine Manipulation \sep Cross-Platform Brigading \sep Algorithmic Governance \sep Artificial Intelligence Regulation
\end{keyword}

\end{frontmatter}

% \linenumbers

% =========================================================================
% SECTION 1: INTRODUCTION
% =========================================================================
\section{Introduction}\label{sec:intro}

Over the past decade, political communication and computational network science have scrutinized the role of encrypted messaging platforms in shaping democratic processes across the Global South \cite{resende2019whatsapp,evangelista2019whatsapp}. In Brazil, India, and across Latin America and Sub-Saharan Africa, WhatsApp established itself as a primary battleground for election campaigns. However, scholarly understanding of this phenomenon has been historically anchored to a rigid architectural constraint: the chat group size limit (originally capped at 256 participants, and later expanded to 1,024) \cite{resende2019whatsapp}. Within that classic paradigm, information dissemination necessarily relied upon horizontal, peer-to-peer forwarding cascades across thousands of fragmented, semi-private groups, creating substantial friction and delay for political organizers.

The 2026 Brazilian Presidential Election represents an unprecedented structural turning point in digital campaign operations. Following Meta's global deployment of \textbf{WhatsApp Channels}---a one-to-many broadcast feature embedded directly into the application's primary navigation under the ``Updates'' tab, operating with no ceiling on subscriber capacity---campaign strategists repurposed the feature from a passive institutional newsletter into an active, synchronized \textbf{Command-and-Control (C2) mobilization weapon}. During the high-stakes runoff between incumbent President Luiz Inácio Lula da Silva (Workers' Party, PT) and Senator Flávio Bolsonaro (Liberal Party, PL), two antagonistic broadcast channels emerged as the operational headquarters of digital warfare:
\begin{enumerate}
    \item \textbf{``EXÉRCITO DO LULA''} (Lula's Army, \texttt{0029VbCbRfD7tkj88k3soZ0Y}): reaching an unprecedented scale of $10.0\text{ million}$ subscribers, managed under direct first-person tactical direction by federal deputy André Janones;
    \item \textbf{``TROPA DO FLÁVIO''} (Flávio's Troop, \texttt{0029Vb8kT6E3AzNNDpBkN22J}): reaching $2.0\text{ million}$ subscribers, operating as a distributed right-wing coalition hub heavily anchored around federal deputy Nikolas Ferreira.
\end{enumerate}

To protect users against viral spam, Meta engineered WhatsApp Channels with a crucial behavioral barrier: push notifications for channel updates are \textbf{silenced by default}. Recognizing that silent delivery destroys real-time synchronization, both campaigns deployed aggressive social-engineering onboarding funnels to bypass this barrier. Channel descriptions and introductory messages explicitly ordered followers to unmute the channel: \textit{``Turn on the notification bell to receive the missions!''} (\textit{Exército do Lula}), while \textit{Tropa do Flávio} disseminated multi-step screen recordings guiding followers through the iOS and Android settings menus to unlock persistent push alerts. By successfully mobilizing millions of citizens to enable push notifications, channel operators gained the ability to dispatch instantaneous, synchronized alerts directly to millions of lock screens simultaneously.

Under Brazilian electoral legislation and updated regulatory directives from Brazil's Superior Electoral Court (TSE Resolution No.\ 23.610), both channels were formally registered in the TSE's public \texttt{DivulgaCandContas} database following court scrutiny, making them officially accountable digital entities. Rather than confining political debates within the chat application, the channels functioned as outbound distribution engines that directed synchronized waves of traffic toward third-party surfaces: Instagram Reels, X (formerly Twitter) Trending Topics, YouTube live streams, and Google Search.

The methodological choice to focus on official WhatsApp Channels rather than mainstream journalistic outlets, digital newspapers, or news channels on YouTube is anchored in the theoretical distinction between \textbf{journalistic mediation} and \textbf{unmediated Command-and-Control (C2)}: news organizations operate with editorial gatekeeping, journalistic guidelines, and independent coverage, whereas WhatsApp Channels represent proprietary, direct-to-voter mobilization infrastructures operated in the first person by campaign tacticians. Within our network modeling framework, WhatsApp Channels function as the \textbf{origin nodes} of coordinated actions, whereas third-party surfaces (such as YouTube streams, Instagram Reels, and Google Search) constitute the \textbf{destination target nodes} where mobilization impacts are observed. To capture the full operational arc of this dynamic while isolating it from standard governmental routines, our empirical inquiry focuses on a rigorous 14-day longitudinal window ($336$ continuous hours from October 17, 2026, 00:00 BRT to October 30, 2026, 23:59 BRT), capturing both the fever pitch of the runoff campaign and the definitive 120-hour post-election aftermath.

This paper addresses three central Research Questions (RQs):
\begin{itemize}
    \item \textbf{RQ1 (Linkless Algorithmic Manipulation):} How do broadcast channels coordinate search engine manipulation without leaving direct hyperlink traces, and how can cross-platform telemetry causally validate this synchronization?
    \item \textbf{RQ2 (Cross-Platform Brigading and Narrative Contestation):} How do opposing political channels mobilize their followers to contest narratives on third-party platforms (such as Instagram and X), and how do their network topologies shape tactical execution?
    \item \textbf{RQ3 (Regulatory Co-Adaptation and Enforcement Resilience):} How do channel operators adapt to proactive judicial interventions by electoral authorities, specifically mandatory artificial intelligence labeling and pre-election suspension threats?
\end{itemize}

The remainder of this article is structured as follows. Section~\ref{sec:related} reviews the literature on coordinated online behavior, search engine manipulation, and algorithmic platform governance. Section~\ref{sec:methodology} formalizes the multi-platform data collection architecture, the event study design, and the tactical taxonomy. Section~\ref{sec:findings} presents structured empirical placeholders for the final post-election results, figures, and statistical analyses. Section~\ref{sec:discussion} outlines policy implications for platform design and electoral regulation, and Section~\ref{sec:conclusion} concludes.

% =========================================================================
% SECTION 2: RELATED WORK AND CONCEPTUAL FRAMEWORK
% =========================================================================
\section{Related Work and Conceptual Framework}\label{sec:related}

Our study intersects three major bodies of literature: coordinated inauthentic and authentic behavior, search engine manipulation, and platform governance under proactive judicial regulation.

\subsection{Coordinated Online Behavior and Outbound Routing}
Research in computational social science has extensively explored Coordinated Link-Sharing Behavior (CLSB) and orchestrated political campaigns on open social networks \cite{giglietto2020coordinated,metaxas2012social}. Giglietto et al. \cite{giglietto2020coordinated} formalized network methods to detect coordinated activity based on the abnormal temporal proximity of URL sharing across Facebook pages and groups. However, while existing coordination models primarily examine link-sharing patterns \textit{within} a single platform or cross-sharing across open web APIs, broadcast channels in encrypted messengers introduce a fundamentally different structural dynamic: \textbf{outbound coordination routing}.

In broadcast channels, coordination does not manifest as horizontal sharing among peers; rather, it originates from a single centralized authority node that broadcasts operational instructions to millions of human agents. Because WhatsApp provides no public comment sections, the operational payoff of these instructions occurs entirely off-platform: driving traffic to an Instagram video to influence its recommendation algorithm, flooding an X hashtag to trigger trending status, or manipulating public search queries.

\subsection{The Search Engine Manipulation Effect (SEME) and Linkless Gaming}
In seminal experimental work, Epstein and Robertson \cite{epstein2015search} introduced the \textbf{Search Engine Manipulation Effect (SEME)}, demonstrating that biased search engine rankings can substantially shift voting preferences among undecided voters without their awareness. While Epstein and Robertson focused on platform-induced algorithmic bias, political actors have long sought to artificially trigger similar shifts via ``Google Bombing''---artificially elevating targeted content for specific search keywords.

Traditionally, Google Bombing required massive networks of webmasters hyperlinking specific anchor texts to a target URL. In our dataset, we identify a contemporary, mobile-first mutation: \textbf{Linkless Search Engine Manipulation}. Channel administrators instruct millions of mobile phone users to manually type identical multi-term search strings into Google simultaneously. Because these commands contain no URLs, they evade automated platform content filters that screen for malicious links or prohibited domains \cite{gillespie2018custodians}. Furthermore, because the resulting queries are generated by millions of genuine, geolocated mobile devices across Brazil, they register as authentic organic interest, tricking search algorithms, trending dashboards, and generative AI overviews into surfacing the campaign's framing.

\subsection{Platform Governance, Algorithmic Auditing, and Judicial Regulation}
The regulation of political discourse on digital platforms has increasingly shifted from voluntary industry self-regulation toward assertive sovereign judicial oversight, particularly in Latin America \cite{gillespie2018custodians,evangelista2019whatsapp}. Brazil's Superior Electoral Court (TSE) stands as a prominent global case study of proactive electoral intervention, exerting judicial authority over platform content, synthetic media, and campaign channels.

Simultaneously, empirical methodologies for \textbf{algorithmic auditing} \cite{hortaribeiro2020auditing} have demonstrated that platform recommendation algorithms frequently exhibit systemic reach disparities across ideological factions. In Brazil's 2026 election, algorithmic audits conducted by civil society coalitions (Sleeping Giants Brasil, Article 19, and CTRL+Z) \cite{sleepinggiants2026nikolas} revealed stark asymmetries in organic video distribution on Instagram. In response to perceived algorithmic throttling on open social feeds, political campaigns increasingly turn to broadcast messaging platforms as an unmediated, chronological communication channel that guarantees direct delivery to voters without algorithmic filtering.

% =========================================================================
% SECTION 3: DATA AND EMPIRICAL METHODOLOGY
% =========================================================================
\section{Data and Empirical Methodology}\label{sec:methodology}

\subsection{Methodological Scope, Platform Selection Rationale, and Eligibility Criteria}

To ensure empirical validity, shield the study against common methodological critiques, and establish transparent boundary conditions, we formalize the operational criteria governing temporal coverage, platform selection, and data inclusion/exclusion:

\paragraph{1. Longitudinal Temporal Delimitation (17--30 October 2026; 336 Continuous Hours):}
Although Brazil's 2026 presidential runoff formally and legally commenced on October 5, 2026 (the day following the first-round ballot on October 4), the empirical observation window consolidated in this study spans exactly \textbf{14 uninterrupted days} (336 continuous hours of data ingestion), commencing at 00:00 BRT on October 17, 2026, and terminating at 23:59 BRT on October 30, 2026. This longitudinal bounding is justified by two substantive political milestones and one crucial computational data integrity parameter:
\begin{itemize}
    \item \textit{The Runoff Campaign Climax (October 17--25, 2026):} Initiated on the eve of the first televised presidential debate of the second round (Rede Bandeirantes, Oct 18)---the precise inflection point when both campaigns fully operationalized their C2 broadcast channels and achieved mass multi-million subscriber mobilization---extending through escalating narrative confrontations, the final televised debate (TV Globo, Oct 23), emergency injunction proceedings before the Superior Electoral Court (TSE), and culminating on Election Day (Sunday, Oct 25);
    \item \textit{The Critical 120-Hour Post-Election Window (October 26--30, 2026):} A 5-day post-election monitoring period is essential to empirically observe victory celebration versus defeat shock, measure subscriber churn and audience attrition, track operational paralysis in the losing coalition, and capture narrative pivoting toward governance or opposition. The definitive cutoff at 23:59 BRT on Friday, October 30, 2026 marks the operational termination of emergency campaign mobilization; extending the window beyond this threshold would introduce confounding noise from routine public administration and bureaucratic appointments unrelated to campaign dynamics;
    \item \textit{Sample Completeness and Algorithmic Buffer Integrity:} Ultra-high-traffic broadcast channels in WhatsApp (exceeding 10 million subscribers and sending dozens of multimedia dispatches daily) experience aggressive client-side DOM element recycling and pagination limits in the WhatsApp Web interface. Anchoring our continuous, exhaustive telemetry from October 17 onward guarantees a 100\% complete census of all broadcast dispatches with zero missing intervals or scrollback truncation artifacts.
\end{itemize}

\paragraph{2. Unit of Analysis and Platform Selection Rationale (Why WhatsApp Channels and Not Journalistic Outlets or YouTube News Channels?):}
A fundamental methodological consideration is why this investigation centers on WhatsApp Channels rather than legacy news media, digital newspapers (e.g., \textit{Folha de S.Paulo}, \textit{g1}, \textit{Metrópoles}), or journalistic television/radio broadcast channels on YouTube (e.g., \textit{Jovem Pan News}, \textit{UOL}, \textit{ICL Notícias}). This decision rests on three network-theoretic and communication arguments:
\begin{enumerate}
    \item \textit{Unmediated Tactical Command (C2) versus Journalistic Editorial Gatekeeping:} Journalistic news organizations and media channels operate with editorial boards, professional fact-checking standards, and journalistic mediation. Our research objective is not to examine \textit{mediated news coverage about the election}, but rather to investigate \textbf{the direct, unmediated algorithmic and tactical Command-and-Control (C2) directives issued by the official campaign commands to their mobilized electorates};
    \item \textit{Two-Mode Graph-Theoretic Asymmetry (Origin vs.\ Target Nodes):} In the bipartite network formalized in this study, WhatsApp Channels serve as the \textbf{origin/emitter nodes ($C$)}, whereas external public platforms (such as YouTube videos, Instagram Reels, X Trending Topics, and Google Search queries) function as the \textbf{destination/target nodes ($T$)}. Treating YouTube news channels or journalistic portals as primary emitters would conflate causes with downstream consequences, because external media outlets are the very surfaces toward which the WhatsApp mobilization armies were directed;
    \item \textit{Deterministic Chronological Delivery versus Algorithmic Curation:} Open social feeds and video platforms deploy proprietary, opaque ranking algorithms that curate content according to commercial engagement metrics. In contrast, WhatsApp Channels deliver broadcasts strictly chronologically directly to subscribers' device notifications, providing an untampered, high-fidelity log of tactical commands as they were generated.
\end{enumerate}

\paragraph{3. Formal Inclusion and Exclusion Criteria:}
\begin{itemize}
    \item \textit{Inclusion of Official Campaign Broadcasts:} 100\% of broadcast messages, external URLs, OpenGraph previews, hashtags, synthetic media tags, and top-4 emoji reaction vectors issued by TSE-registered official runoff campaign accounts during the observation window were captured and analyzed;
    \item \textit{Exclusion of Private Conversational Groups:} Traditional WhatsApp chat groups (limited to 256 or 1,024 members) were deliberately excluded. First, capturing private conversational groups without explicit user consent risks harvesting personal telephone numbers, profile names, and private discussions, violating data protection legislation (specifically Brazil's General Data Protection Law, LGPD, Law No.\ 13.709/2018). Second, conversational groups exhibit decentralized horizontal forwarding cascades whose sampling is heavily skewed by the private invite links accessible to researchers. By contrast, official WhatsApp Channels are completely public broadcasts, contain zero citizen telephone numbers or PII, and permit an exhaustive census ($100\%$ sample) of the campaign's broadcast history;
    \item \textit{Exclusion of Apocryphal or Unofficial Channels:} Unverified grassroots channels or speculative accounts not formally registered in the TSE's \texttt{DivulgaCandContas} database were excluded to guarantee internal causal validity and ensure that all observed dispatches represent official, authorized campaign strategy;
    \item \textit{Exclusion of Peripheral Grassroots Social Initiatives:} Spontaneous post-election socialization initiatives emerging among supporters (such as grassroots dating/networking chats) were excluded from formal network modeling, as they constitute playful community leisure rather than coordinated political command-and-control operations.
\end{itemize}

\subsection{Multi-Channel Longitudinal Ingestion Architecture}
Our primary data collection pipeline ingested all broadcast messages from the official runoff campaign channels registered with Brazil's Superior Electoral Court (TSE) under Resolution No.\ 23.610:
\begin{itemize}
    \item \textbf{EXÉRCITO DO LULA} (\texttt{0029VbCbRfD7tkj88k3soZ0Y}): $10.0\text{ million}$ subscribers;
    \item \textbf{TROPA DO FLÁVIO} (\texttt{0029Vb8kT6E3AzNNDpBkN22J}): $2.0\text{ million}$ subscribers;
    \item Allied satellite channels tracked for cross-channel spillover: \textit{Canal Oficial do Lula} (\texttt{0029VaDH1GU7DAWpJTrCXH39}) and \textit{Nikolas Ferreira - Oficial} (\texttt{0029VbDNtfo4o7qVYbwemy1f}).
\end{itemize}

We developed a headless browser DOM extraction pipeline (\texttt{whatsapp\_channel\_extractor.js}) that connects to an active WhatsApp Web session and executes programmatic backward scrolling to capture the complete longitudinal stream. For each message $m_i$, the extractor captures:
\begin{equation}
m_i = \left\langle \text{msgId}, t_{\text{UTC}}, t_{\text{BRT}}, \text{Text}, \mathcal{U}_{\text{links}}, \mathcal{H}_{\text{tags}}, \mathcal{M}_{\text{media}}, \mathcal{R}_{\text{emojis}} \right\rangle
\end{equation}
where $\mathcal{U}_{\text{links}}$ denotes extracted external URLs and OpenGraph preview metadata, $\mathcal{H}_{\text{tags}}$ contains all explicit hashtags, $\mathcal{M}_{\text{media}}$ identifies media attachments and regulatory badges (such as native ``AI Content'' / \textit{Conteúdo de IA} labels), and $\mathcal{R}_{\text{emojis}}$ records the ordered vector of top-4 reaction emojis with corresponding aggregated counts.

\subsection{Triangulation with External Telemetry Streams}
To evaluate the cross-platform impact of internal WhatsApp dispatches, we triangulated channel logs against four external telemetry data sources:
\begin{enumerate}
    \item \textbf{Google Trends Relative Search Volume (RSV):} Hourly search interest queries for Brazil ($0$--$100$ index) retrieved via programmatic endpoints, capturing normalized query volumes for targeted search bombing terms;
    \item \textbf{X (Twitter) Brazil Trending Topics:} Hourly snapshots of the Top-50 Trending Topics in Brazil, parsed to record exact ranks, durations, and syntactic keyword matches;
    \item \textbf{Wikimedia REST API Pageviews:} Daily and hourly absolute desktop and mobile pageviews for Portuguese Wikipedia biographical articles directly tied to campaign search bombing queries;
    \item \textbf{Public Channel Subscriber Probing:} Unauthenticated hourly HTTP probes targeting public WhatsApp channel invite pages to record longitudinal subscriber counts $S_c(t)$ and measure post-election churn.
\end{enumerate}

\subsection{Causal Event Study Framework}
To establish causal linkages between a WhatsApp dispatch at timestamp $t_0$ and external platform reactions, we implement a quasi-experimental event study design. For an observed time window $t \in [t_0 - \Delta, t_0 + \Delta]$ (where $\Delta$ ranges from 6 to 24 hours depending on the platform's telemetry granularity), we define the baseline pre-treatment activity $\bar{Y}_{\text{pre}}$ and measure the post-treatment shock magnitude $\delta$:
\begin{equation}
\delta = \frac{\max_{t \in [t_0, t_0 + \Delta]} Y(t) - \bar{Y}_{\text{pre}}}{\sigma(Y_{\text{pre}}) + \epsilon}
\end{equation}
where $Y(t)$ represents the external telemetry metric (e.g., Google Trends RSV, Wikipedia daily views, or X trending rank). A causal shock is validated when: (i) $\bar{Y}_{\text{pre}} \approx 0$; (ii) the spike occurs within $\tau \le 60\text{ minutes}$ of the WhatsApp dispatch; and (iii) the syntactic fingerprint in external data matches the exact phrasing broadcast on the channel.

\subsection{Tactical Mission Taxonomy}
We categorized each dispatch into an 8-fold tactical taxonomy:
\begin{enumerate}
    \item \textbf{Internal Command \& Troop Morale:} Dispatches delivering narrative framing, ideological reinforcement, and spiritual or emotional appeals without outbound links;
    \item \textbf{X/Twitter Hashtag Battles:} Direct directives commanding subscribers to tweet and retweet specific hashtags to capture the Trending Topics;
    \item \textbf{Cross-Platform Content Distribution:} Routine dissemination of official candidate speeches, institutional announcements, and campaign clips hosted on Instagram, Facebook, or web portals;
    \item \textbf{Cross-Platform Brigading:} Direct calls to flood specific third-party posts (Instagram Reels, YouTube streams) with coordinated likes, comments, and reports;
    \item \textbf{Linkless Google Search Bombing:} Explicit commands ordering followers to manually type targeted search queries into Google;
    \item \textbf{Synchronized Event Mobilization:} Real-time crisis coordination during live televised debates and electoral rallies;
    \item \textbf{Gamified Recruitment \& Notification Push:} Directive dispatches instructing users to unmute the channel bell and forward invite links to personal contact lists;
    \item \textbf{Synthetic Media with Regulatory AI Labeling:} Dispatches sharing AI-generated images or videos adhering to TSE disclosure rules.
\end{enumerate}

\subsection{Ethical Framework and Data Governance}
All data analyzed in this study were collected from public broadcast channels formally registered with Brazil's electoral authorities. The dataset contains no personally identifiable information (PII), phone numbers, private chat messages, or user-level identifiers. Our protocol complies strictly with the Brazilian General Data Protection Law (LGPD, Law No.\ 13.709/2018), the ethical guidelines for social media research established by the Association of Internet Researchers (AoIR), and the regulatory parameters established by TSE Resolution No.\ 23.610.

% =========================================================================
% SECTION 4: EMPIRICAL FINDINGS (STRUCTURED PLACEHOLDERS)
% =========================================================================
\section{Empirical Findings}\label{sec:findings}

% [PLACEHOLDER NOTE: This entire empirical section will be populated with the complete
% longitudinal dataset once the 120-hour post-election window closes on Friday, October 30, 2026.]

\subsection{RQ1: Linkless Google Search Bombing and Multi-Platform Causal Shocks}

% [PLACEHOLDER TABLE 1: SUMMARY OF RECORDED DATASET METRICS]
\begin{table}[htb]
\centering
\caption{Comprehensive Longitudinal Dataset Metrics ($N=1,331$ Messages across the 14-day Observation Window, 17/10 to 30/10/2026).}
\label{tab:metrics_dataset}
\begin{tabular}{@{}lcc@{}}
\toprule
\textbf{Analytical Dimension} & \textbf{EXÉRCITO DO LULA} & \textbf{TROPA DO FLÁVIO} \\ \midrule
Channel ID & \texttt{0029VbCbRfD7tkj88k3soZ0Y} & \texttt{0029Vb8kT6E3AzNNDpBkN22J} \\
Audited Peak Followers & $10.0\text{ million}$ & $2.1\text{ million}$ \\
Post-Election Retention (30/10) & $10.0\text{ million}$ & $1.9\text{ million (-200k churn / -9.5\%)}$ \\
Total Messages Processed ($N$) & $1,182$ & $149$ \\
Outbound External Links Dispatched & $160$ & $41$ \\
Coordinated X Hashtag Directives & $123$ ($134$ missions) & $3$ ($4$ missions) \\
Google Search Bombing Campaigns & $4\text{ explicit waves (7 dispatches)}$ & $0\text{ waves}$ \\
Primary Destination Platform & Instagram ($109$), WA ($14$), X ($12$) & Instagram ($32$), X ($3$), YT ($2$) \\
AI Content Badges Detected & $0$ & $3$ \\ \bottomrule
\end{tabular}
\end{table}

% [PLACEHOLDER FIGURE 1: CAUSAL SHOCK TIMELINE (THE TALVANE ALBUQUERQUE CASE)]
\begin{figure}[htb]
\centering
\fbox{\parbox{0.92\textwidth}{\centering \vspace{3.8cm} \textbf{[PLACEHOLDER FIGURE 1: Multi-Platform Causal Shock Timeline --- Talvane Albuquerque Case]}\\
High-resolution four-tier event study diagram illustrating the exact causal chain: (a) WhatsApp command at 15:36 BRT ($8.1\text{M}$ recipients); (b) Google Trends Relative Search Volume surge from $0 \rightarrow 72 \rightarrow 100$ in 60 minutes; (c) X Trending Topics entry at \#1 with simultaneous imperative verb match (\texttt{"Joguem"}); and (d) Portuguese Wikipedia pageview explosion ($+139,266\%$) and second-hop article traversal to Ceci Cunha. \vspace{3.8cm}}}
\caption{Triangulated multi-platform causal shock following a linkless Google Search Bombing command on WhatsApp.}
\label{fig:causal_shock_placeholder}
\end{figure}

\noindent \textit{[PLACEHOLDER ANALYSIS RQ1: In-depth statistical analysis of the 4 waves of Google Search Bombing, including econometric event-study estimations, decay half-lives, and evidence of AI Overview suppression by Google algorithms.]}

\subsection{RQ2: Cross-Platform Brigading and the Narrative Conflict over Northeast Brazil}

% [PLACEHOLDER FIGURE 2: LONGITUDINAL STREAMGRAPH OF TOPICAL TRAJECTORIES]
\begin{figure}[htb]
\centering
\fbox{\parbox{0.92\textwidth}{\centering \vspace{3.8cm} \textbf{[PLACEHOLDER FIGURE 2: Definitive 14-Day Longitudinal Streamgraph (17/10 to 30/10/2026)]}\\
ThemeRiver / Streamgraph visualization mapping the dynamic flow of tactical missions and narrative controversies over time, featuring key inflection points: Debate da Band (18/10), Talvane Albuquerque Shock (22/10), Debate da Globo and TSE Blackout Threat (23/10), Election Day Crisis Operations (25/10), and Post-Election Divergence (26/10 to 30/10). \vspace{3.8cm}}}
\caption{Dynamic Streamflow of Coordinated Tactical Categories across the 2026 Runoff and Post-Election Window.}
\label{fig:streamgraph_placeholder}
\end{figure}

\noindent \textit{[PLACEHOLDER ANALYSIS RQ2: Detailed comparative evaluation of the cross-platform clash surrounding deputy Nikolas Ferreira's muted Instagram Reel, the copyright censorship narrative, and the asymmetric retaliation executed by Lula's army.]}

\subsection{RQ3: Judicial Enforcement, AI Disclosures, and Pre-Election Blackout Gamification}

% [PLACEHOLDER TABLE 2: TSE REGULATORY ACTIONS AND CAMPAIGN RESPONSES]
\begin{table}[htb]
\centering
\caption{[PLACEHOLDER] Chronicle of Superior Electoral Court (TSE) Judicial Actions and Empirical Campaign Co-Adaptation.}
\label{tab:tse_placeholder}
\begin{tabular}{@{}p{3.5cm}p{3.5cm}p{6.0cm}@{}}
\toprule
\textbf{TSE Regulatory Event} & \textbf{Timestamp in Logs} & \textbf{Observed Campaign Co-Adaptation} \\ \midrule
Mandatory Channel Registration & Oct 17--19, 2026 & Both channels registered in \texttt{DivulgaCandContas} to avert deactivation. \\
Mandatory AI Labeling Enforcement & Oct 19, 2026 & \textit{Tropa do Flávio} applies native \texttt{"Conteúdo de IA"} labels to 3 synthetic assets. \\
48-Hour Channel Blackout Threat & Oct 23, 2026 & \textit{Exército do Lula} gamifies suspension threat into a scarcity drive ($+1.9\text{M}$ followers in 28h). \\
Post-Election Paralysis \& Defeat & Oct 25--30, 2026 & \textit{Tropa do Flávio} freezes dispatches; subscriber churn of $-100,000$ followers recorded. \\ \bottomrule
\end{tabular}
\end{table}

% [PLACEHOLDER FIGURE 3: POST-ELECTION RETENTION AND SUBSCRIBER CHURN]
\begin{figure}[htb]
\centering
\fbox{\parbox{0.92\textwidth}{\centering \vspace{3.8cm} \textbf{[PLACEHOLDER FIGURE 3: Post-Election Divergence and Subscriber Churn Curves]}\\
Longitudinal curves contrasting the post-election trajectory: \textit{Exército do Lula} maintaining 10M subscribers and expanding secondary channels versus \textit{Tropa do Flávio} experiencing immediate command paralysis after 25/10 18:22 BRT and dropping from 2.1M to 2.0M followers. \vspace{3.8cm}}}
\caption{Empirical Post-Election Divergence: Command Freeze and Audience Churn.}
\label{fig:churn_placeholder}
\end{figure}

\noindent \textit{[PLACEHOLDER ANALYSIS RQ3: Quantitative analysis of the gamified scarcity mechanism deployed during the TSE blackout debate, accompanied by an audit of synthetic media labeling compliance.]}

% =========================================================================
% SECTION 5: DISCUSSION AND POLICY IMPLICATIONS (PLACEHOLDER)
% =========================================================================
\section{Discussion and Policy Implications}\label{sec:discussion}

\noindent \textit{[PLACEHOLDER DISCUSSION: Detailed scholarly discussion addressing the systemic vulnerabilities of broadcast messaging platforms in electoral contexts. Key topics will include: (i) how push-notification engineering circumvents platform friction; (ii) the challenge of regulating linkless algorithmic manipulation under existing election laws; (iii) the unintended consequences of publicized regulatory blackout threats; and (iv) design recommendations for platform architectures to maintain accountability without undermining end-to-end user privacy.]}

% =========================================================================
% SECTION 6: LIMITATIONS AND FUTURE WORK (PLACEHOLDER)
% =========================================================================
\section{Limitations and Future Work}\label{sec:limitations}

\noindent \textit{[PLACEHOLDER LIMITATIONS: Discussion of methodological boundaries, including the lack of internal exposure telemetry (impressions per subscriber) within WhatsApp Channels, reliance on external public proxies, and potential sampling biases. Proposed future avenues: multi-modal video/audio semantic processing and cross-national comparative studies across upcoming elections in other Global South democracies.]}

% =========================================================================
% SECTION 7: CONCLUSION (PLACEHOLDER)
% =========================================================================
\section{Conclusion}\label{sec:conclusion}

\noindent \textit{[PLACEHOLDER CONCLUSION: Definitive concluding synthesis highlighting the transition from conversational chat groups to mass broadcast C2 hubs as the central socio-technical paradigm of 2026 political mobilization, reinforcing the empirical validity of linkless search manipulation and multi-platform telemetry triangulation.]}

% =========================================================================
% BACKMATTER
% =========================================================================
\section*{Declaration of Competing Interest}
The authors declare that they have no known competing financial interests or personal relationships that could have appeared to influence the work reported in this paper.

\section*{Funding}
Open Access publication of this article is supported by the Coordenação de Aperfeiçoamento de Pessoal de Nível Superior -- Brasil (CAPES) under the CAPES--Elsevier Transformative Agreement (Portaria CAPES No.\ 120/2024).

\section*{Data and Code Availability}
In accordance with Open Science and FAIR (Findable, Accessible, Interoperable, Reusable) data principles, all anonymized message datasets ($N=1,331$ messages), variable codebooks, bipartite network projections in GEXF format, and Python extraction scripts are publicly available in the official research repository on GitHub at: \url{https://github.com/alkawamo/brasnam-eleicoes-2026}. Furthermore, an interactive, bilingual replication pipeline is provided via Google Colab (\url{https://colab.research.google.com/github/alkawamo/brasnam-eleicoes-2026/blob/main/BRASNAM_2026_Colab_Pipeline.ipynb}), enabling zero-install 1-click cloud execution and visual replication of all analytical figures and event study models.

\bibliographystyle{elsarticle-num}
\bibliography{references_osnem}

\end{document}
